\section{Knapsack 0/1}
\whenuse{Elegir cada objeto como m\'aximo una vez, con capacidad moderada, para
maximizar valor o decidir qu\'e sumas son alcanzables.}

Cada objeto puede elegirse a lo sumo una vez; por eso la capacidad se recorre
en orden descendente.

\begin{lstlisting}[style=cpp]
ll knapsack_01(const vector<int> &weight, const vector<ll> &value,
               int capacity) {
    assert(weight.size() == value.size() && capacity >= 0);
    assert(all_of(weight.begin(), weight.end(), [](int w) { return w >= 0; }));
    vector<ll> dp(capacity + 1, 0);
    for (int item = 0; item < (int)weight.size(); ++item)
        for (int current = capacity; current >= weight[item]; --current)
            dp[current] = max(dp[current],
                              dp[current - weight[item]] + value[item]);
    return *max_element(dp.begin(), dp.end());
}
\end{lstlisting}

\complexity{$O(nW)$ tiempo y $O(W)$ memoria.}

\section{Bitmask DP: camino hamiltoniano m\'inimo}
\whenuse{El estado depende del subconjunto exacto de elementos ya usados y
$n$ es peque\~no, normalmente $n\le 20$.}

Encuentra el costo m\'inimo de un camino que visita cada v\'ertice exactamente
una vez, con inicio y fin libres; no exige volver al inicio. La matriz cuadrada
\texttt{cost[u][v]} representa un grafo dirigido: usar \texttt{LINF} si no existe
la arista $u\to v$. Para un grafo no dirigido, cargar ambos sentidos.
Admite costos negativos y devuelve \texttt{LINF} si no existe un camino completo;
con un solo v\'ertice devuelve cero. Devuelve el costo, no la secuencia de nodos.
Se precalculan m\'ascaras de aristas existentes para recorrer s\'olo los vecinos
que a\'un no fueron visitados.

\begin{lstlisting}[style=cpp]
ll hamiltonian_path(const vector<vector<ll>> &cost) {
    int n = cost.size();
    assert(1 <= n && n < 31);
    assert(all_of(cost.begin(), cost.end(), [n](const auto &row) {
        return (int)row.size() == n;
    }));
    int states = 1 << n;
    vector<int> outgoing(n, 0);
    for (int u = 0; u < n; ++u)
        for (int v = 0; v < n; ++v)
            if (cost[u][v] != LINF) outgoing[u] |= 1 << v;
    vector<vector<ll>> dp(states, vector<ll>(n, LINF));
    for (int v = 0; v < n; ++v) dp[1 << v][v] = 0;
    for (int mask = 1; mask < states; ++mask)
        for (int last = 0; last < n; ++last) if (dp[mask][last] != LINF)
            for (int nexts = outgoing[last] & ~mask; nexts; nexts &= nexts - 1) {
                int next = __builtin_ctz((unsigned)nexts); // nexts != 0
                dp[mask | (1 << next)][next] = min(
                    dp[mask | (1 << next)][next],
                    dp[mask][last] + cost[last][next]);
            }
    return *min_element(dp.back().begin(), dp.back().end());
}
\end{lstlisting}

\complexity{$O(2^n n^2)$ tiempo y $O(2^n n)$ memoria; usar para $n$ peque\~no.}

\warning{Requiere $1\le n<31$ para la m\'ascara, pero tiempo y memoria suelen
limitarlo mucho antes: con $n=20$, s\'olo los valores de la DP ocupan 160 MiB,
m\'as el costo de los vectores. Los costos finitos y todas las sumas parciales
deben estar estrictamente entre $-\texttt{LINF}$ y $\texttt{LINF}$.
Una cota suficiente, para $n>1$, es $(n-1)\max|w|<\texttt{LINF}$ sobre las
aristas existentes. Nunca sumar la marca de arista inexistente.}
